\documentclass[../../main.tex]{subfiles}
\graphicspath{{../../image/}{../../figure/}}
\begin{document}

\section{比较定理与 $L^1$ 稳定性}\label{section:ch11-比较定理与L1稳定性}

\cref{chapter:ch06-标量守恒律L1理论}的断言将在此处以更精细的形式, 针对一维真正非线性标量守恒律\eqref{align:ch11-11-1-1}的特殊情形重新建立。关键因素是函数

\begin{align}
Q(u,v,w)=
\begin{cases}
f(v)-f(u)-\dfrac{f(u)-f(w)}{u-w}[v-u], & \text{若 }u\neq w,\\[3mm]
f(v)-f(u)-f'(u)[v-u], & \text{若 }u=w,
\end{cases}\label{align:ch11-11-8-1}
\end{align}
对 $u,v,w\in\mathbb{R}$ 定义。显然, $Q(u,v,w)=Q(w,v,u)$。由于 $f$ 一致凸, 当 $v$ 位于 $u$ 与 $w$ 之间时 $Q(u,v,w)$ 为负, 当 $v$ 位于端点为 $u$ 与 $w$ 的区间之外时为正。特别地, 对 Burgers 方程(4.2.1), $Q(u,v,w)=\dfrac12(v-u)(v-w)$。

第一步是精细化序性质。

\begin{theorem}\label{theorem:ch11-11-8-1}
设 $u$ 与 $\bar u$ 是上半平面上\eqref{align:ch11-11-1-1}的可容许解, 对应初值 $u_0$ 与 $\bar u_0$, 满足
\begin{align}
u_0(x)\leqslant\bar u_0(x),\quad \text{对所有 }x\in(y,\bar y).\label{align:ch11-11-8-2}
\end{align}
设 $\psi(\cdot)$ 是与解 $u$ 关联、从点 $(y,0)$ 出发的任意前向特征线, $\bar\psi(\cdot)$ 是与 $\bar u$ 关联、从点 $(\bar y,0)$ 出发的任意前向特征线。则对任意 $t>0$ 且 $\psi(t)<\bar\psi(t)$,
\begin{align}
u(x,t)\leqslant\bar u(x,t),\quad \text{对所有 }x\in(\psi(t),\bar\psi(t)).\label{align:ch11-11-8-3}
\end{align}
\end{theorem}

\begin{proof}
固定任意区间 $(z,\bar z)$, $\psi(t)<z<\bar z<\bar\psi(t)$。考虑与解 $u$ 关联、从点 $(z,t)$ 出发的极大后向特征线 $\xi(\cdot)$, 以及与 $\bar u$ 关联、从点 $(\bar z,t)$ 出发的极小后向特征线 $\bar\xi(\cdot)$。于是 $\xi(0)\geqslant y$, $\bar\xi(0)\leqslant\bar y$。

先假设 $\xi(0)<\bar\xi(0)$。积分方程
\begin{align}
\partial_t[u-\bar u]+\partial_x[f(u)-f(\bar u)]=0\label{align:ch11-11-8-4}
\end{align}
在梯形 $\{(x,\tau):0<\tau<t,\ \xi(\tau)<x<\bar\xi(\tau)\}$ 上, 应用 Green 公式得
\begin{align}
&\int_z^{\bar z}[u(x,t)-\bar u(x,t)]\,\mathrm{d}x-\int_{\xi(0)}^{\bar\xi(0)}[u_0(x)-\bar u_0(x)]\,\mathrm{d}x\nonumber\\
&\quad=-\int_0^t Q(u(\xi(\tau),\tau),\bar u(\xi(\tau),\tau),u(\xi(\tau),\tau))\,\mathrm{d}\tau\nonumber\\
&\qquad\quad-\int_0^t Q(\bar u(\bar\xi(\tau),\tau),u(\bar\xi(\tau),\tau),\bar u(\bar\xi(\tau),\tau))\,\mathrm{d}\tau.\label{align:ch11-11-8-5}
\end{align}
\eqref{align:ch11-11-8-5}右端两个积分均非正。因此由\eqref{align:ch11-11-8-2}, $u(\cdot,t)-\bar u(\cdot,t)$ 在 $(z,\bar z)$ 上的积分非正。

现假设 $\xi(0)\geqslant\bar\xi(0)$。则直线 $\xi(\cdot)$ 与 $\bar\xi(\cdot)$ 必在某时刻 $s\in[0,t]$ 相交。此时在三角形 $\{(x,\tau):s<\tau<t,\ \xi(\tau)<x<\bar\xi(\tau)\}$ 上积分方程\eqref{align:ch11-11-8-4}, 用同样论证推出 $u(\cdot,t)-\bar u(\cdot,t)$ 在 $(z,\bar z)$ 上的积分仍非正。

由于 $(z,\bar z)$ 是 $(\psi(t),\bar\psi(t))$ 的任意子区间, 我们得出\eqref{align:ch11-11-8-3}。证毕。
\end{proof}

作为上述定理的推论, 我们推出函数 $u(\cdot,t)-\bar u(\cdot,t)$ 在 $(-\infty,\infty)$ 上的变号数随时间不增。事实上, 假设有诸点 $-\infty=y_0<y_1<\cdots<y_n<y_{n+1}=\infty$, 使得在每个区间 $(y_i,y_{i+1})$ 上, $u_0(\cdot)-\bar u_0(\cdot)$ 当 $i$ 偶时非负, 当 $i$ 奇时非正。设 $\psi_i(\cdot)$ 是与解 $u$ 关联、从点 $(y_i,0)$ 出发的任意前向特征线, $i$ 奇; $\bar\psi_i(\cdot)$ 是与 $\bar u$ 关联、从点 $(y_i,0)$ 出发的任意前向特征线, $i$ 偶。这些曲线一般按如下规则被赋予有限寿命: 在最早碰撞时刻 $t_1$, 某条 $\psi_i$ 与某条 $\bar\psi_j$ 碰撞, 这两条曲线终止。然后在下一次碰撞时刻 $t_2$, 任意(存活的)$\psi_k$ 与 $\bar\psi_\ell$ 碰撞, 这两条曲线同样终止; 依此类推。由\cref{theorem:ch11-11-8-1}, 对任意 $t\in[0,t_1]$, $u(\cdot,t)-\bar u(\cdot,t)$ 经历 $n$ 次变号; 对任意 $t\in[t_1,t_2]$, 经历 $n-2$ 次变号; 依此类推。特别地, 所谓\textbf{回折数}(lap number), 即解 $u(\cdot,t)$ 的图像与任意固定常数 $\bar u$ 的交叉数, 随时间不增。

由\cref{chapter:ch06-标量守恒律L1理论}, 标量守恒律任意一对可容许解的空间 $L^1$ 距离随时间不增。在当前设定下, 我们将证明可确定 $L^1$ 距离严格递减的条件及递减速率。


\begin{theorem}\label{theorem:ch11-11-8-2}
设 $u$ 与 $\bar u$ 是\eqref{align:ch11-11-1-1}的可容许解, 初值 $u_0,\bar u_0\in L^1(-\infty,\infty)$。则 $\|u(\cdot,t)-\bar u(\cdot,t)\|_{L^1(-\infty,\infty)}$ 是 $t$ 的不增函数, 在 $(0,\infty)$ 上局部 Lipschitz。对任意固定 $t\in(0,\infty)$, 考虑(可能空且至多可数)集
\begin{align}
\begin{cases}
\mathscr{J}=\{y\in(-\infty,\infty):u_+<\bar u_+\leqslant\bar u_-<u_-\},\\
\bar{\mathscr{J}}=\{y\in(-\infty,\infty):\bar u_+<u_+\leqslant u_-<\bar u_-\},
\end{cases}\label{align:ch11-11-8-6}
\end{align}
其中 $u_\pm,\bar u_\pm$ 分别表示 $u(y\pm,t),\bar u(y\pm,t)$。令
\begin{align}
u_*=
\begin{cases}
u_\pm, & \text{若 }u_+=u_-,\\
u_-, & \text{若 }u_+<u_-\text{ 且 }\dfrac{f(u_+)-f(u_-)}{u_+-u_-}\geqslant\dfrac{f(\bar u_+)-f(\bar u_-)}{\bar u_+-\bar u_-},\\
u_+, & \text{若 }u_+<u_-\text{ 且 }\dfrac{f(u_+)-f(u_-)}{u_+-u_-}<\dfrac{f(\bar u_+)-f(\bar u_-)}{\bar u_+-\bar u_-},
\end{cases}\label{align:ch11-11-8-7a}
\end{align}
\begin{align}
\bar u_*=
\begin{cases}
\bar u_\pm, & \text{若 }\bar u_+=\bar u_-,\\
\bar u_-, & \text{若 }\bar u_+<\bar u_-\text{ 且 }\dfrac{f(\bar u_+)-f(\bar u_-)}{\bar u_+-\bar u_-}\geqslant\dfrac{f(u_+)-f(u_-)}{u_+-u_-},\\
\bar u_+, & \text{若 }\bar u_+<\bar u_-\text{ 且 }\dfrac{f(\bar u_+)-f(\bar u_-)}{\bar u_+-\bar u_-}<\dfrac{f(u_+)-f(u_-)}{u_+-u_-}.
\end{cases}\label{align:ch11-11-8-7b}
\end{align}
则
\begin{align}
\frac{d^+}{dt}\|u(\cdot,t)-\bar u(\cdot,t)\|_{L^1(-\infty,\infty)}
=2\sum_{y\in\mathscr{J}}Q(u_-,\bar u_*,u_+)+2\sum_{y\in\bar{\mathscr{J}}}Q(\bar u_-,u_*,\bar u_+).\label{align:ch11-11-8-8}
\end{align}
\end{theorem}

\begin{proof}
首先在 $u(\cdot,t)-\bar u(\cdot,t)$ 在 $(-\infty,\infty)$ 上经历有限次变号的特殊情形建立\eqref{align:ch11-11-8-8}。于是假设有诸点 $-\infty=y_0<y_1<\cdots<y_n<y_{n+1}=\infty$, 使得 $u(\cdot,t)-\bar u(\cdot,t)$ 在区间 $(y_i,y_{i+1})$ 上当 $i$ 偶时非负, 当 $i$ 奇时非正。特别地, 任意 $y\in\mathscr{J}$ 必是某 $y_i$, $i$ 奇; 任意 $y\in\bar{\mathscr{J}}$ 必是某 $y_i$, $i$ 偶。

设 $\psi_i(\cdot)$ 是与解 $u$ 关联、从点 $(y_i,t)$ 出发的(唯一)前向特征线, $i$ 奇; $\bar\psi_i(\cdot)$ 是与 $\bar u$ 关联、从点 $(y_i,t)$ 出发的前向特征线, $i$ 偶。固定 $s>t$, $s-t$ 足够小, 使得在 $[t,s]$ 上上述曲线无碰撞。在区域 $\{(x,\tau):t<\tau<s,\ \psi_i(\tau)<x<\bar\psi_{i+1}(\tau)\}$, $i$ 奇, 以及区域 $\{(x,\tau):t<\tau<s,\ \bar\psi_i(\tau)<x<\psi_{i+1}(\tau)\}$, $i$ 偶, 上积分\eqref{align:ch11-11-8-4}。应用 Green 公式并利用\cref{theorem:ch11-11-8-1}, 推出
\begin{align}
&\|u(\cdot,s)-\bar u(\cdot,s)\|_{L^1}-\|u(\cdot,t)-\bar u(\cdot,t)\|_{L^1}\nonumber\\
&\quad=\sum_{i\text{ 偶}}\int_{\bar\psi_i(s)}^{\bar\psi_{i+1}(s)}[u(x,s)-\bar u(x,s)]\,\mathrm{d}x
+\sum_{i\text{ 奇}}\int_{\psi_i(s)}^{\psi_{i+1}(s)}[\bar u(x,s)-u(x,s)]\,\mathrm{d}x\nonumber\\
&\qquad-\sum_{i\text{ 偶}}\int_{y_i}^{y_{i+1}}[u(x,t)-\bar u(x,t)]\,\mathrm{d}x
-\sum_{i\text{ 奇}}\int_{y_i}^{y_{i+1}}[\bar u(x,t)-u(x,t)]\,\mathrm{d}x\nonumber\\
&\quad=\sum_{i\text{ 奇}}\int_t^s\{Q(u(\psi_i(\tau)-,\tau),\bar u(\psi_i(\tau)-,\tau),u(\psi_i(\tau)+,\tau))\nonumber\\
&\qquad\qquad\qquad+Q(u(\psi_i(\tau)+,\tau),\bar u(\psi_i(\tau)+,\tau),u(\psi_i(\tau)-,\tau))\}\,\mathrm{d}\tau\nonumber\\
&\qquad+\sum_{i\text{ 偶}}\int_t^s\{Q(\bar u(\bar\psi_i(\tau)-,\tau),u(\bar\psi_i(\tau)-,\tau),\bar u(\bar\psi_i(\tau)+,\tau))\nonumber\\
&\qquad\qquad\qquad+Q(\bar u(\bar\psi_i(\tau)+,\tau),u(\bar\psi_i(\tau)+,\tau),\bar u(\bar\psi_i(\tau)-,\tau))\}\,\mathrm{d}\tau.\label{align:ch11-11-8-9}
\end{align}
由\cref{theorem:ch11-11-3-1}, 当 $s\downarrow t$ 时, \eqref{align:ch11-11-8-9}右端第一个积分中的被积函数, 若 $y_i\notin\mathscr{J}$ 则趋于零, 若 $y_i\in\mathscr{J}$ 则趋于 $2Q(u_-,\bar u_*,u_+)$。类似地, 第二个积分中的被积函数, 若 $y_i\notin\bar{\mathscr{J}}$ 则趋于零, 若 $y_i\in\bar{\mathscr{J}}$ 则趋于 $2Q(\bar u_-,u_*,\bar u_+)$。因此, 将\eqref{align:ch11-11-8-9}除以 $s-t$ 并令 $s\downarrow t$, 即得\eqref{align:ch11-11-8-8}。

现在转向一般情形, 即 $u(\cdot,t)-\bar u(\cdot,t)$ 在 $(-\infty,\infty)$ 上可能经历无穷次变号, 注意此时开集 $\{x\in(-\infty,\infty):u(x\pm,t)-\bar u(x\pm,t)<0\}$ 是不相交开区间 $(y_i,\bar y_i)$ 的可数并。对 $m=1,2,\cdots$, 设 $u_m$ 是\eqref{align:ch11-11-1-1}在 $(-\infty,\infty)\times[t,\infty)$ 上的可容许解, 初值
\begin{align}
u_m(x,t)=
\begin{cases}
\bar u(x,t), & x\in\bigcup\limits_{i=m}^\infty(y_i,\bar y_i),\\
u(x,t), & \text{其他}.
\end{cases}\label{align:ch11-11-8-10}
\end{align}
于是 $u_m(\cdot,t)-\bar u(\cdot,t)$ 在 $(-\infty,\infty)$ 上经历有限次变号, 因而对 $\tau\geqslant t$, $\dfrac{d^+}{d\tau}\|u_m(\cdot,\tau)-\bar u(\cdot,\tau)\|_{L^1}$ 由\eqref{align:ch11-11-8-8}的类似式给出。此外, 函数 $\tau\mapsto\dfrac{d^+}{d\tau}\|u_m(\cdot,\tau)-\bar u(\cdot,\tau)\|_{L^1}$ 在 $t$ 处右连续, 且右连续模与 $m$ 无关。为验证此点, 注意小跳对 $\|u_m(\cdot)-\bar u(\cdot)\|_{L^1}$ 变化率的总贡献很小, 由 $u(\cdot,t)$ 与 $\bar u(\cdot,t)$ 在 $(-\infty,\infty)$ 上的全变差控制, 而(有限个)大跳的贡献由\cref{theorem:ch11-11-3-1}右连续。因此, 令 $m\to\infty$ 取极限, 我们对一般解 $u$ 与 $\bar u$ 建立了\eqref{align:ch11-11-8-8}。证毕。
\end{proof}

由上述定理, $u(\cdot,t)$ 与 $\bar u(\cdot,t)$ 的 $L^1$ 距离仅当其中一个函数的图像在跳跃间断点处与另一个函数的图像交叉时才可能减小。更鲁棒的收缩由一种新度量实现, 它用专为两个解定制的权来加权 $L^1$ 距离。

对 $v,\bar v\in BV(-\infty,\infty)$, 令
\begin{align}
\rho(v,\bar v)=\int_{-\infty}^\infty\bigl\{(V(x)+\bar V(\infty)-\bar V(x))[v(x)-\bar v(x)]^+\nonumber\\
\qquad+(\bar V(x)+V(\infty)-V(x))[\bar v(x)-v(x)]^+\bigr\}\,\mathrm{d}x,\label{align:ch11-11-8-11}
\end{align}
其中上标 $+$ 表示"正部", $w^+=\max\{w,0\}$, $V$ 或 $\bar V$ 表示 $v$ 或 $\bar v$ 的变差函数, $V(x)=TV_{(-\infty,x)}v(\cdot)$, $\bar V(x)=TV_{(-\infty,x)}\bar v(\cdot)$。

\begin{theorem}\label{theorem:ch11-11-8-3}
设 $u$ 与 $\bar u$ 是\eqref{align:ch11-11-1-1}的可容许解, 初值 $u_0,\bar u_0\in BV(-\infty,\infty)$。则对任意固定 $t\in(0,\infty)$,
\begin{align}
\frac{d^+}{dt}\rho(u(\cdot,t),\bar u(\cdot,t))
\leqslant&-\int_{-\infty}^\infty Q(u(x,t),\bar u(x,t),u(x,t))\,\mathrm{d}V_t^c(x)\nonumber\\
&-\int_{-\infty}^\infty Q(\bar u(x,t),u(x,t),\bar u(x,t))\,\mathrm{d}\bar V_t^c(x)\nonumber\\
&-\sum_{y\in\mathscr{K}}(u_--u_+)Q(u_-,\bar u_*,u_+)
-\sum_{y\in\bar{\mathscr{K}}}(\bar u_--\bar u_+)Q(\bar u_-,u_*,\bar u_+)\nonumber\\
&+(V_t(\infty)+\bar V_t(\infty))\left\{\sum_{y\in\mathscr{J}}Q(u_-,\bar u_*,u_+)
+\sum_{y\in\bar{\mathscr{J}}}Q(\bar u_-,u_*,\bar u_+)\right\},\label{align:ch11-11-8-12}
\end{align}
其中 $V_t$ 或 $\bar V_t$ 是 $u(\cdot,t)$ 或 $\bar u(\cdot,t)$ 的变差函数; $V_t^c$ 或 $\bar V_t^c$ 表示 $V_t$ 或 $\bar V_t$ 的连续部分; $u_\pm,\bar u_\pm$ 表示 $u(y\pm,t),\bar u(y\pm,t)$; $u_*,\bar u_*$ 仍由\eqref{align:ch11-11-8-7a}与\eqref{align:ch11-11-8-7b}确定; 集 $\mathscr{J},\bar{\mathscr{J}}$ 由\eqref{align:ch11-11-8-6}定义, $\mathscr{K}$ 或 $\bar{\mathscr{K}}$ 表示 $u(\cdot,t)$ 或 $\bar u(\cdot,t)$ 的跳跃点集:
\begin{align}
\begin{cases}
\mathscr{K}=\{y\in(-\infty,\infty):u_+<u_-\},\\
\bar{\mathscr{K}}=\{y\in(-\infty,\infty):\bar u_+<\bar u_-\}.
\end{cases}\label{align:ch11-11-8-13}
\end{align}
\end{theorem}

\begin{proof}
如\cref{theorem:ch11-11-8-2}的证明那样开始: 假设实轴有分划 $-\infty=y_0<y_1<\cdots<y_n<y_{n+1}=\infty$, 使得在每个区间 $(y_i,y_{i+1})$ 上, $u(\cdot,t)-\bar u(\cdot,t)$ 当 $i$ 偶时非负, 当 $i$ 奇时非正。考虑与 $u$ 关联、从每个 $y_i,t$ 出发的前向特征线 $\psi_i(\cdot)$, $i$ 奇; 以及与 $\bar u$ 关联、从每个 $y_i,t$ 出发的前向特征线 $\bar\psi_i(\cdot)$, $i$ 偶。

我们聚焦于某个 $(y_i,y_{i+1})$, $i$ 偶。仅讨论 $-\infty<y_i<y_{i+1}<\infty$ 的情形, 其他情形更简单。除 $\bar\psi_i(\cdot)$ 外, 下文考虑的所有特征线均与解 $u$ 关联。论证随前向特征线 $\chi_0$ 从 $(y_i,t)$ 出发位于 $\bar\psi_i(\cdot)$ 左侧还是右侧而有所不同; 为确定起见, 讨论后一种略复杂的情形。

固定小正数 $\varepsilon$, 识别所有 $z_1,\cdots,z_N$, $y_i<z_1<\cdots<z_N<y_{i+1}$, 使得 $u(z_I-,t)-u(z_I+,t)\geqslant\varepsilon$, $I=1,\cdots,N$。考虑从点 $(z_I,t)$ 出发的前向特征线 $\chi_I(\cdot)$, $I=1,\cdots,N$。然后选 $s>t$, $s-t$ 足够小, 使得:
\begin{enumerate}[(a)]
\item 在时间区间 $[t,s]$ 上, 特征线 $\chi_0,\chi_1,\cdots,\chi_N,\psi_{i+1}$ 之间无相交;
\item 对 $I=1,\cdots,N$, 若 $\zeta_I(\cdot)$ 与 $\xi_I(\cdot)$ 表示从点 $(\chi_I(s),s)$ 出发的极小与极大后向特征线, 则 $u(\cdot,t)$ 在区间 $(\zeta_I(t),z_I)$ 与 $(z_I,\xi_I(t))$ 上的全变差不超过 $\varepsilon/N$;
\item 若 $\zeta(\cdot)$ 表示从 $(\psi_{i+1}(s),s)$ 出发的极小后向特征线, 则 $u(\cdot,t)$ 在区间 $(\zeta(t),y_{i+1})$ 上的全变差不超过 $\varepsilon$;
\item 若 $\zeta_0(\cdot)$ 是从 $(\psi_i(s),s)$ 出发的极小后向特征线, $\xi_0(\cdot)$ 是从 $(\chi_0(s),s)$ 出发的极大后向特征线, 则 $u(\cdot,t)$ 在区间 $(\zeta_0(t),y_i)$ 与 $(y_i,\xi_0(t))$ 上的全变差不超过 $\varepsilon$。
\end{enumerate}

对 $I=0,\cdots,N-1$ 和稍后固定的 $k$, 在区间 $[\chi_I(s),\chi_{I+1}(s)]$ 上取分划: $\chi_I(s)=x_I^0<x_I^1<\cdots<x_I^k<x_I^{k+1}=\chi_{I+1}(s)$; 类似地在 $[\chi_N(s),\psi_{i+1}(s)]$ 上: $\chi_N(s)=x_N^0<x_N^1<\cdots<x_N^k<x_N^{k+1}=\psi_{i+1}(s)$。对 $I=0,\cdots,N$, $j=1,\cdots,k$, 考虑从点 $(x_I^j,s)$ 出发的极大后向特征线 $\zeta_I^j(\cdot)$, 识别其在 $t$ 时刻的截距 $z_I^j=\zeta_I^j(t)$。并设 $z_0^0=y_i$, $z_N^{k+1}=y_{i+1}$, $z_{I-1}^{k+1}=z_I^0=z_I$, $I=1,\cdots,N$。

现注意恒等式
\begin{align}
R-S=-D,\label{align:ch11-11-8-14}
\end{align}
其中
\begin{align}
R=&\int_{\bar\psi_i(s)}^{\chi_0(s)}V_t(y_i)[u(x,s)-\bar u(x,s)]\,\mathrm{d}x\nonumber\\
&+\sum_{I=0}^N\sum_{j=0}^k\int_{x_I^j}^{x_I^{j+1}}V_t(z_I^j+)[u(x,s)-\bar u(x,s)]\,\mathrm{d}x,\label{align:ch11-11-8-15}\\
S=&\sum_{I=0}^N\sum_{j=0}^k\int_{z_I^j}^{z_I^{j+1}}V_t(z_I^j+)[u(x,t)-\bar u(x,t)]\,\mathrm{d}x,\label{align:ch11-11-8-16}\\
D=&\sum_{I=0}^N\sum_{j=1}^k\int_t^s[V_t(z_I^j+)-V_t(z_I^{j-1}+)]Q(u(\zeta_I^j(\tau),\tau),\bar u(\zeta_I^j(\tau)-,\tau),u(\zeta_I^j(\tau),\tau))\,\mathrm{d}\tau\nonumber\\
&+\sum_{I=1}^N\int_t^s[V_t(z_I+)-V_t(z_{I-1}+)]Q(u(\chi_I(\tau)-,\tau),\bar u(\chi_I(\tau)-,\tau),u(\chi_I(\tau)+,\tau))\,\mathrm{d}\tau\nonumber\\
&+\int_t^s[V_t(y_i+)-V_t(y_i)]Q(u(\chi_0(\tau)-,\tau),\bar u(\chi_0(\tau)-,\tau),u(\chi_0(\tau)+,\tau))\,\mathrm{d}\tau\nonumber\\
&-\int_t^s V_t(y_i)Q(\bar u(\bar\psi_i(\tau)-,\tau),u(\bar\psi_i(\tau)+,\tau),\bar u(\bar\psi_i(\tau)+,\tau))\,\mathrm{d}\tau\nonumber\\
&-\int_t^s V_t(z_N^k+)Q(u(\psi_{i+1}(\tau)-,\tau),\bar u(\psi_{i+1}(\tau)-,\tau),u(\psi_{i+1}(\tau)+,\tau))\,\mathrm{d}\tau.\label{align:ch11-11-8-17}
\end{align}
为验证\eqref{align:ch11-11-8-14}, 先在区域 $\{(x,\tau):t<\tau<s,\ \bar\psi_i(\tau)<x<\chi_0(\tau)\}$, $\{(x,\tau):t<\tau<s,\ \zeta_I^j(\tau)<x<\zeta_I^{j+1}(\tau)\}$, $\{(x,\tau):t<\tau<s,\ \chi_I(\tau)<x<\xi_I^1(\tau)\}$, $\{(x,\tau):t<\tau<s,\ \xi_I^k(\tau)<x<\chi_{I+1}(\tau)\}$, $\{(x,\tau):t<\tau<s,\ \xi_N^k(\tau)<x<\psi_{i+1}(\tau)\}$ 上积分\eqref{align:ch11-11-8-4} 并应用 Green 公式; 然后以权重 $V_t(y_i),V_t(z_I^j+),V_t(z_I+),V_t(z_N^k+)$ 对所得方程作加权和。

为估计 $R$, 注意 $V_t(y_i)\geqslant V_s(\chi_0(s))$, 且 $V_t(z_I^j+)\geqslant V_s(x_I^j+)$, $I=0,\cdots,N$, $j=0,\cdots,k$。因此, 若取 $x_I^{j+1}-x_I^j$ 充分小, 可保证
\begin{align}
R\geqslant\int_{\bar\psi_i(s)}^{\psi_{i+1}(s)}V_s(x)[u(x,s)-\bar u(x,s)]\,\mathrm{d}x-(s-t)\varepsilon.\label{align:ch11-11-8-18}
\end{align}

为估计 $S$, 只需注意 $V_t(\cdot)$ 非减, 故
\begin{align}
S\leqslant\int_{\bar\psi_i(t)}^{\psi_{i+1}(t)}V_t(x)[u(x,t)-\bar u(x,t)]\,\mathrm{d}x.\label{align:ch11-11-8-19}
\end{align}

为估计 $D$, 注意由 $Q$ 的性质, 五项均非负。对 $I=0,\cdots,N$, $j=1,\cdots,k$, $V_t(z_I^j+)-V_t(z_I^{j-1}+)\geqslant V_t^c(z_I^j)-V_t^c(z_I^{j-1})$。此外,
\begin{align}
Q(u(\zeta_I^j(\tau),\tau),\bar u(\zeta_I^j(\tau)-,\tau),u(\zeta_I^j(\tau),\tau))=Q(u(z_I^j,t),\bar u(p_\tau(z_I^j),t),u(z_I^j,t)),\label{align:ch11-11-8-20}
\end{align}
其中单调递增函数 $p_\tau$ 由
\begin{align}
p_\tau(x)=x+(\tau-t)[f'(u(x,t))-f'(\bar u(p_\tau(x),t))]\label{align:ch11-11-8-21}
\end{align}
确定。取 $x_I^{j+1}-x_I^j$ 充分小, 使 $V_t^c(\cdot)$ 在每个区间 $(z_I^j,z_I^{j+1})$ 上的振荡不超过 $\varepsilon$, Stieltjes 积分的标准估计蕴含
\begin{align}
\sum_{I=0}^N\sum_{j=1}^k[V_t(z_I^j+)-V_t(z_I^{j-1}+)]Q(u(\zeta_I^j(\tau),\tau),\bar u(\zeta_I^j(\tau)-,\tau),u(\zeta_I^j(\tau),\tau))\nonumber\\
\geqslant\int_{y_i}^{y_{i+1}}Q(u(x,t),\bar u(p_\tau(x),t),u(x,t))\,\mathrm{d}V_t^c(x)-c\varepsilon.\label{align:ch11-11-8-22}
\end{align}

现在合并\eqref{align:ch11-11-8-14}与\eqref{align:ch11-11-8-18}、\eqref{align:ch11-11-8-19}、\eqref{align:ch11-11-8-17}、\eqref{align:ch11-11-8-22}, 然后将所得不等式除以 $s-t$, 令 $s\downarrow t$, 最后令 $\varepsilon\downarrow0$。这给出
\begin{align}
&\frac{d^+}{dt}\int_{\bar\psi_i(t)}^{\psi_{i+1}(t)}V_t(x)[u(x,t)-\bar u(x,t)]\,\mathrm{d}x\nonumber\\
&\quad\leqslant-\int_{y_i}^{y_{i+1}}Q(u(x,t),\bar u(x,t),u(x,t))\,\mathrm{d}V_t^c(x)\nonumber\\
&\qquad-\sum(u_--u_+)Q(u_-,\bar u_*,u_+)+V_t(y_i)Q(\bar u_-,u_*,\bar u_+)\nonumber\\
&\qquad+V_t(y_{i+1})Q(u_-,\bar u_*,u_+),\label{align:ch11-11-8-23}
\end{align}
其中求和遍历 $\mathscr{K}\cap(y_i,y_{i+1})$ 中所有 $y$, 以及 $y_i$ 若 $y_i\in\mathscr{K}$ 且 $\chi_0$ 位于 $\bar\psi_i$ 右侧。$u_\pm,\bar u_\pm,u_*,\bar u_*$ 当然在相应 $y$ 处取值。

接下来我们聚焦于区间 $(y_i,y_{i+1})$, $i$ 奇。完全对称的论证给出, 代替\eqref{align:ch11-11-8-23},
\begin{align}
&\frac{d^+}{dt}\int_{\psi_i(t)}^{\bar\psi_{i+1}(t)}(V_t(\infty)-V_t(x))[\bar u(x,t)-u(x,t)]\,\mathrm{d}x\nonumber\\
&\quad\leqslant-\int_{y_i}^{y_{i+1}}Q(u(x,t),\bar u(x,t),u(x,t))\,\mathrm{d}V_t^c(x)\nonumber\\
&\qquad-\sum(u_--u_+)Q(u_-,\bar u_*,u_+)+(V_t(\infty)-V_t(y_i+))Q(u_-,\bar u_*,u_+)\nonumber\\
&\qquad+(V_t(\infty)-V_t(y_{i+1}))Q(\bar u_-,u_*,\bar u_+),\label{align:ch11-11-8-24}
\end{align}
其中求和遍历 $\mathscr{K}\cap(y_i,y_{i+1})$ 中所有 $y$, 以及 $y_{i+1}$ 若 $y_{i+1}\in\mathscr{K}$ 且与 $u$ 关联、从点 $(y_{i+1},t)$ 出发的前向特征线位于 $\bar\psi_{i+1}(\cdot)$ 左侧。

于是对所有偶 $i$ 写出\eqref{align:ch11-11-8-23}, 对所有奇 $i$ 写出\eqref{align:ch11-11-8-24}, 并对 $i=0,\cdots,n$ 求和。这给出
\begin{align}
&\frac{d^+}{dt}\int_{-\infty}^\infty\bigl\{V_t(x)[u(x,t)-\bar u(x,t)]^+
+(V_t(\infty)-V_t(x))[\bar u(x,t)-u(x,t)]^+\bigr\}\,\mathrm{d}x\nonumber\\
&\quad\leqslant-\int_{-\infty}^\infty Q(u(x,t),\bar u(x,t),u(x,t))\,\mathrm{d}V_t^c(x)\nonumber\\
&\qquad-\sum_{y\in\mathscr{K}}(u_--u_+)Q(u_-,\bar u_*,u_+)
+V_t(\infty)\left\{\sum_{y\in\mathscr{J}}Q(u_-,\bar u_*,u_+)
+\sum_{y\in\bar{\mathscr{J}}}Q(\bar u_-,u_*,\bar u_+)\right\}.\label{align:ch11-11-8-25}
\end{align}

通过如\cref{theorem:ch11-11-8-2}证明中的技术论证, 可证明即使允许 $u(\cdot,t)-\bar u(\cdot,t)$ 在 $(-\infty,\infty)$ 上经历无穷次变号, \eqref{align:ch11-11-8-25}仍成立。

交换 $u$ 与 $\bar u$ 的角色, 写出由\eqref{align:ch11-11-8-25}得到的不等式, 再与\eqref{align:ch11-11-8-25}合并。这给出\eqref{align:ch11-11-8-12}。证毕。
\end{proof}

估计\eqref{align:ch11-11-8-12}是精细的, 至少对分段光滑解它以等式成立。\eqref{align:ch11-11-8-12}右端所有项除 $-\sum_{y\in\mathscr{J}}(u_--u_+)Q(u_-,\bar u_*,u_+)$ 与 $-\sum_{y\in\bar{\mathscr{J}}}(\bar u_--\bar u_+)Q(\bar u_-,u_*,\bar u_+)$ 外均为负。然而即使这些正项也被负项 $V_t(\infty)Q(u_-,\bar u_*,u_+)$ 与 $\bar V_t(\infty)Q(\bar u_-,u_*,\bar u_+)$ 抵消。因此距离函数 $\rho(u(\cdot,t),\bar u(\cdot,t))$ 一般严格递减。

泛函 $\rho$ 的类似物将在\cref{chapter:ch14-前沿追踪与标准riemann半群}中用于建立方程组守恒律解的 $L^1$ 稳定性。


\end{document}
